% Part: first-order-logic
% Chapter: beyond
% Section: many-sorted-logic

\documentclass[../../../include/open-logic-section]{subfiles}

\begin{document}

\olfileid[tr]{fol}{byd}{msl}

\olsection{Çok Sıralı Mantık}

Birinci dereceden mantıkta değişkenler ile niceleyiciler tek bir !!{domain}
üzerinde değer alır. Oysa birden çok (ayrık) !!{domain}s kullanmak çoğu zaman
yararlıdır: örneğin sayılar için bir !!a{domain}, geometrik nesneler için bir
!!a{domain}, sayılardan sayılara fonksiyonlar için bir !!a{domain}, değişmeli
gruplar için bir !!a{domain} vb. isteyebiliriz.

Çok sıralı mantık bu tür bir çerçeve sağlar. Önce bir ``sıralar'' listesiyle
başlanır; bir nesnenin ``sırası'', içinde yer alması öngörülen ``!!{domain}''i
belirtir. Ardından her sıra için !!{variable}s ve niceleyiciler, ayrıca
(genellikle) her sıra için bir !!{identity} bulunur. Fonksiyonlar ile
bağıntılar da argüman olarak alabilecekleri nesnelerin sıralarıyla
``tiplendirilir.'' Bunun dışında birinci dereceden mantığın olağan kuralları
korunur; niceleyici kurallarının bir sürümü her sıra için yinelenir.

Örneğin uluslararası ilişkileri incelemek için iki nesne sırası, Fransız
yurttaşları ile Alman yurttaşları, içeren bir dil seçebiliriz. İlk sıradaki
nesneler için ``şarap içer'' tekli bağıntısı; ikinci sıradaki nesneler için
``sosis yer'' başka bir tekli bağıntısı; ilk argümanı ilk, ikinci argümanı da
ikinci sıradan olan iki argüman alan ``çok uluslu evli çift oluşturur'' ikili
bağıntısı bulunabilir. $a$, $b$, $c$ değişkenlerini Fransız yurttaşları;
$x$, $y$, $z$ değişkenlerini Alman yurttaşları üzerinde değer alacak biçimde
kullanırsak
\[
\lforall[a][\lforall x][(\Atom{\Obj{MarriedTo}}{a,x} \lif
(\Atom{\Obj{DrinksWine}}{a} \lor \lnot \Atom{\Obj{EatsWurst}}{x})]]
\]
ifadesi, herhangi bir Fransız bir Almanla evliyse ya Fransızın şarap içtiğini
ya da Almanın sosis yemediğini bildirir.

Çok sıralı mantık doğal bir biçimde birinci dereceden mantığa gömülebilir.
Bunun için çok sıralı !!{domain}s içindeki bütün nesneler tek bir birinci
dereceden !!{domain} içinde bir araya getirilir, sıraları izlemek için tekli
!!{predicate}s kullanılır ve niceleyiciler görelileştirilir. Örneğin yukarıdaki
örneğe karşılık gelen birinci dereceden dil, betimlenen diğer bağıntılara ek
olarak ``$\Obj{German}$'' ve ``$\Obj{French}$'' tekli !!{predicate}s{}ini içerir;
sıra koşulları ise silinir. Sıralı bir $\lforall[x][!A]$ niceleyicisi, burada
$x$ Alman sırasından bir !!a{variable} olmak üzere, şu ifadeye çevrilir:
\[
\lforall[x][(\Atom{\Obj{German}}{x} \lif !A)].
\]
Sıraların ayrık olmasını güvence altına alan aksiyomlar---örneğin
$\lforall[x][\lnot (\Atom{\Obj{German}}{x} \land
  \Atom{\Obj{French}}{x})]$---ve ``şarap içer'' bağıntısının yalnızca
$\Atom{\Obj{French}}{x}$ yüklemini sağlayan nesneler için geçerli olmasını
güvence altına alan aksiyomlar vb. eklememiz gerekir. Bu uzlaşımlar ve
aksiyomlarla çok sıralı !!{sentence}s{}in birinci dereceden !!{sentence}s{}e,
çok sıralı !!{derivation}s{}in da birinci dereceden !!{derivation}s{}e çevrildiğini
göstermek zor değildir. Çok sıralı !!{structure}s da karşılık gelen birinci
dereceden !!{structure}s{}a ve tersine ``çevrilir''; dolayısıyla çok sıralı
mantık için de bir tamlık teoremimiz vardır.

\end{document}
